\chapter{Algorithmic Design and Implementation}
\label{c:algorithmic-impl}

%%%%%%%%%%%%%%%%%%%%%%%%%%%%%%%%%%%%%%%%%%%%%%%%%%%%%%%%%%%%%%%%%%%%%%%%%%%%%%%%

% ==========================================
% Section: Core Data Structures
% ==========================================
\section{Core Data Structures}
To construct the foundation of the simulator, we must first define the core data structures that map the mathematical framework into efficient computational memory.

% ==========================================
% Subsection: Defining the State
% ==========================================
\subsection{Defining the State}
In quantum mechanics a state is defined as a linear combination of orthogonal basis states and theirs corresponding complex probability amplitude.
So to be able to store a state we need to find a structure that can hold this data and link the amplitude to the correct basis state.
We define the basis state as a bitstring corresponding to the occupation number representation of the Fock space \eqref{eq:occupation_number_representation}, allowing it to be stored as an unsigned integer.
Its corresponding probability amplitude can be stored as a complex number.
\cite[Ch. 1.2.1]{sakurai2020modern}
Because of the fact, that we have a dynamic number of basis states that can also change in the runtime a lot of static structures, like dense arrays, are not feasible.
Rather than relaying on static approaches like a dense array, which would scale exponentially and store unnecessary zero-amplitude states, we implement a sparse representation using a hash map.
This library provides a more efficient, cache friendly and faster unordered hash map which suits our requirements.
\cite{ankerl_unordered_dense}


Combining these architectural choices, the quantum state is defined via a type alias:
\begin{minted}{cpp}
    using State = ankerl::unordered_dense::map<uint64_t, std::complex<double>>;
\end{minted}

% ==========================================
% Subsection: Defining the Generators and Hamiltonian
% Todo: Reference?
% ==========================================
\subsection{Defining the Generators and Hamiltonian}
Both, excitation generators \eqref{eq:mult_fe_exc_generator} and Hamiltonian rely on the second-quantized operators, which means they can get stored with the same structure.
The \texttt{FermionTerm} structure encodes creation and annihilation orbital information using integer vectors.
The complex scalar \texttt{weight} gets included to support arbitrary fermionic operators.
\begin{minted}{cpp}
    struct FermionTerm {
        std::vector<int> creation_idx;
        std::vector<int> annihilation_idx;
        std::complex<double> weight = 1.0;
    };
\end{minted}
In order to build combinations the FermionTerm structure can also be stored in a vector.
For readability and a convenient way to input Googles OpenFermion string representation can be used as in- and output.
\cite[chapter II.A]{mcclean2019openfermion}

% ==========================================
% Section: Matrix-Free Direct State Evolution
% ==========================================
\section{Matrix-Free Direct State Evolution}\label{sec:matrix-free-direct-state-evolution}
The core computational challenge in simulating many-body quantum systems lies in the efficient application of excitation operators to the underlying state vector. Constructing and applying global unitary matrices is highly memory-prohibitive for large Hilbert spaces. Therefore, the simulator leverages a direct state-vector approach, exploiting the algebraic properties of the excitation generator to execute updates directly on the computational basis.

To translate this theoretical framework into a high-performance algorithmic routine, we must deconstruct the evaluation of the unitary operation. As an initial step, we isolate the fundamental mechanics by considering the qubit excitation framework. This isolated model provides a computational baseline that is then subsequently generalized to incorporate fermionic anticommutation relations.


% ==========================================
% Subsection: Fermionic Excitation
% ==========================================
\subsection{Fermionic Excitation}\label{subsec:fermionic-excitation}
Rather than relying on full matrix exponentiation, the direct state evolution approach exploits the algebraic properties of the generator to update the state.
We can simplify the evaluation of the unitary operation, in order to compute the excitation efficiently.

Consider an initial quantum state on which the excitation is applied.
Recalling the generalized generator from Equation \eqref{eq:mult_fe_exc_generator}, the Generator is fully defined by the orbital sets $\boldsymbol{p}$ and $\boldsymbol{q}$, storing which electrons are exchanged.
The unitary itself also requires an angle $\theta$, which we can for now ignore.
With these information, the Unitary posses all the necessary parameters, and it can be expanded as:
\begin{equation} \label{eq:fermionic-excitation-term}
    U(\theta) = e^{-i\frac{\theta}{2}G^{\boldsymbol{k}}_{\boldsymbol{l}}} = \cos\left(\frac{\theta}{2}\right) I - i\sin\left(\frac{\theta}{2}\right)G^{\boldsymbol{k}}_{\boldsymbol{l}} + \left(1 - \cos\left(\frac{\theta}{2}\right)\right) P_0.
\end{equation}
Here, $P_0$ represents the projector onto the nullspace of the generator.
This nullspace is spanned by all computational basis states $\ket{\boldsymbol{i}}$ for which $G^{\boldsymbol{k}}_{\boldsymbol{l}} \ket{\boldsymbol{i}} = 0$.

Furthermore, when applying a specific configuration $\ket{\Phi}$, the unitary acts in a trivial way (as the identity operator) if the configuration is contained within the nullspace of the generator.
Otherwise, it induces a rotation between the original configuration and the $n$-fold excited excitation:
\begin{equation} \label{eq:qubit-excitation-unitary-transformation-configuration-differences}
    U(\theta) \ket{\Phi} = \begin{cases} \ket{\Phi}, & \text{if } P_0\ket{\Phi} = \ket{\Phi} \\ \left(\cos\left(\frac{\theta}{2}\right) I - i\sin\left(\frac{\theta}{2}\right)G\right)\ket{\Phi}, & \text{else.} \end{cases}
\end{equation}
\cite[chapter II.A]{kottmann2020feasible}

To expand this generalized view to be able to handle fermions we need to have a deeper understanding of the creation $\hat{a}^{\dagger}_k$ and annihilation $\hat{a}_k$ operators.
The fermionic anticommutation defined in \eqref{sec:the-second-quantization-framework} induces a parity phase, when acting on a target spin-orbital $k$, determined by the total occupation of all preceding orbitals in the canonical ordering.
Algebraically, the action of these operators on Fock state $\ket{\boldsymbol{n}}$ is expressed as:
\begin{align} \label{eq:annhiliation-creation-algebraic-definition}
    a_k \ket{\mathbf{n}} &= \delta_{n_k, 1} (-1)^{\sum_{j=0}^{k-1} n_j} \ket{n_0, \dots, n_k - 1, \dots, n_{M-1}}, \\
    a_k^\dagger \ket{\mathbf{n}} &= \delta_{n_k, 0} (-1)^{\sum_{j=0}^{k-1} n_j} \ket{n_0, \dots, n_k + 1, \dots, n_{M-1}},
\end{align}~\cite[Box 1.1]{helgaker2000molecular}
where $\delta$ represents the Kronecker delta, enforcing the Pauli exclusion principle defined in \eqref{sec:fermionic-many-body-systems}.
The phase factor $(-1)^{\sum_{j=0}^{k-1} n_j}$ evaluates to $+1$ for an even occupation count and $-1$ for an odd count.

% Todo: To be revised
To evaluate an $n$-fold excitation, constituent operators are applied sequentially, with the occupation value updated immediately after each step ($n_k \mapsto 1 - n_k$). This algebraic tracking of the state vector naturally calculates the required anticommutation phases, completely eliminating the need for complex Jordan-Wigner strings.

This sequential update enforces the strict antisymmetry required by the Pauli exclusion principle. Because the algorithm relies purely on initial and target orbital occupancies, the accurate evaluation of the system depends entirely on correctly applying this classically computed parity phase.

Finally, this approach offers a significant computational advantage over physical quantum hardware. By determining the phase through direct occupation counting, the simulator bypasses the severe gate overhead and noise caused by the deep CNOT circuits and Pauli-Z chains required on actual quantum devices.
% end: to be revised
\cite[chapter III.B]{anand2022quantum}\cite[chapter 2.1.1]{schleich2025quantum}

% Implementation (new chapter?)
In order to implement such complex behaviors we need to have a deeper understanding for all the components and formulate an algorithm, which computes the excitation.
Therefore, we examine the two rotational components of the term \eqref{eq:qubit-excitation-unitary-transformation-configuration-differences}.
The first part governed by $\cos\left(\frac{\theta}{2}\right) I$, expresses a direct scaling of the initial state's probability amplitude, which is computationally straightforward and can be expressed as:
\begin{equation}
    c_{\boldsymbol{i}} \leftarrow c_{\boldsymbol{i}} \cos\left(\frac{\theta}{2}\right).\label{eq:qubit-excitation-update-basis-state}
\end{equation}

The second component involves the coupling and is more complex as it involves the action of the generator.
The term $-i\sin\left(\frac{\theta}{2}\right)G^{\boldsymbol{k}}_{\boldsymbol{l}}$ can be decomposed into two operational steps.
First, the algorithm computes the partner state $\ket{\boldsymbol{i'}}$ which can be computed by applying the generator on the basis state.
\begin{equation}
    c \ket{\boldsymbol{i'}} = -iG^{\boldsymbol{k}}_{\boldsymbol{l}} \ket{\boldsymbol{i}}, \label{eq:qubit-excitation-baseline-partner-state}
\end{equation}~\cite[Appendix - Eigenstates of fermionic generators]{kottmann2020feasible}
The derivation of this equation can be read in~\ref{subsec:derivation---calculate-fermionic-ecitation-partner-state}.
Here, the partner state represents the electronic configuration generated by the excitation.
The transitional coefficient $c$ is restricted to real values and even $c = \pm 1$.
The sign of $c$ depends on the direction of the particle exchange (i.e. whether the operation acts as an excitation or de-excitation)

Finally, this term induces an amplitude update on the newly acquired partner state.
The corresponding probability amplitude is updated as:
\begin{equation}
	c_{\boldsymbol{i'}} \leftarrow c_{\boldsymbol{i'}} + c \sin\left(\frac{\theta}{2}\right). \label{eq:qubit-excitation-partner-state-update-phase}
\end{equation}
\cite[chapter II.A]{kottmann2020feasible}

% --- Implementation ---
The C++ implementation follows the explained algorithm.
A brief explanation of the function would look like this:
\begin{minted}[frame=lines, framesep=2mm, obeytabs=true, tabsize=4]{text}
Input: state, fermionicTerm, theta

For each basis_state in state:
    If basis_state is in Nullspace:
        Keep basis_state unchanged
    Else:
        Determine operator direction
        Apply annihilation and creation (track fermionic phase)
        Generate partner_state
        Update amplitudes of basis_state and partner_state

Return new_state
\end{minted}

To ensure our efficiency goals we made the following architectural decisions.
\begin{itemize}
    \item The memory for the new state gets preallocated based on the input states size.
    \item To evaluate the nullspace projector we bypass normal conditional logic and instead use bit-shifting and masking directly on the state.
    \item The fermionic parity phase (which requires counting the number of proceeding electrons) is evaluated using the Cpu's native population count.
    Therefore, avoiding unnecessary counting loops.
    \item Partner basis states get generated using XOR bitwise operation rather than initializing intermediate state objects.
    \item Orbital sets get passed as references to avoid copying standard integer vectors.
\end{itemize}

The completed C++ function is in the appendix under~\ref{subsec:appendix-fermionic-excitation}
% ==========================================
% Subsection: Fermionic SWAP (fSWAP)
% ==========================================
\subsection{Fermionic SWAP (fSWAP)}
Apart from applying excitations, the simulation of many-body fermionic systems often requires reordering fermionic orbitals to satisfy hardware constraints or algorithmic requirements.
This reordering must preserve fermionic anticommutation relations and is achieved via the fermionic SWAP operation.
In the computational basis, this gate is defined as:
\begin{equation}
	U_{\text{fSWAP}} = \begin{pmatrix}
						   1 & 0 & 0 & 0 \\
						   0 & 0 & 1 & 0 \\
						   0 & 1 & 0 & 0 \\
						   0 & 0 & 0 & -1
	\end{pmatrix}
\end{equation}
The negative phase on the $\ket{11}$ element enforces the required exchange sign when swapping two occupied orbitals.
\cite[chapter II.]{Verstraete2009quantum}

% --- Implementation ---

As already seen before we also here avoid constructing dense operator matrices and therefore evaluate the fSWAP gate directly on the state vector.
Furthermore, iterating exclusively over the non-zero amplitudes reduces time and memory complexity.
The C++ code for the function is also part of the appendix (\ref{subsec:appendix-fswap}).


% ==========================================
% Subsection: Expectation Value Evaluation
% ==========================================
\subsection{Expectation Value Evaluation}\label{sec:expectation-value-evaluation}
Extracting physical observables requires evaluating expectation values $\bra{\phi} H \ket{\psi}$, which is called matrix element if $\phi \neq \psi$.

Decomposing the Hamiltonian into its linear combination of fermionic operator strings $H = \sum_k w_k \hat{O}_k$, and expanding the states to its computational basis ($\ket{\psi} = \sum_x c_x \ket{x}$ and $\bra{\phi} = \sum_y d_y^* \bra{y}$), yields us that for each operator term $\hat{O}_k$ applies:
\begin{equation}
    \bra{\phi} \hat{O}_k \ket{\psi} = \sum_{x} \sum_{y} d_y^* c_x \bra{y} \hat{O}_k \ket{x}
\end{equation}
Note that, applying an operator term $\hat{O}_k$ to a basis state $\ket{x}$ either destroys the state, if the operation violates the Pauli exclusion principle, or maps it into a new basis state $\ket{x'_k}$ with a parity phase $p_{k,x} \in \{-1, 1\}$.
Or expressed mathematically $\hat{O}_k \ket{x} = p_{k,x} \ket{x'_k}$.
Due to $\bra{y}$ and $\ket{x}$ have orthonormal bases, the inner product evaluates to a Kronecker delta:
\begin{equation}
    \bra{y} \ket{x'_k} = \delta_{y, x'_k} = \begin{cases} 1 & \text{if } y = x'_k \\ 0 & \text{if } y \neq x'_k \end{cases}
\end{equation}
This collapses the double summation into a linear sum over the non-zero amplitudes of $\ket{\psi}$.
The total expectation value can therefore be computed as:
\begin{equation}
    \bra{\phi} H \ket{\psi} = \sum_k w_k \sum_{x} d_{x'_k}^* p_{k,x} c_x\label{eq:expectation-value-decomposed}
\end{equation}
Algorithmically, resolving $\ket{x'_k}$ via bitwise manipulations allows the simulator to evaluate observables without constructing dense matrix operators.
\cite[chapter II.A, II.B]{kohda2022quantum}\cite[Chapter 1.2.1, 1.3.3, 1.4.3, 7.7]{sakurai2020modern}\cite[chapter 6.6]{benito_kottmann_qc_2025}

% --- Implementation ---

Following this decomposition, the simulator evaluates the expectation value by breaking the Hamiltonian down and accumulating each Hamiltonian term by itself.
For each fermionic operator string $\hat{O}_k$, the algorithm iterates only over the non-zero amplitudes of $\ket{\psi}$.
Annihilation and creation operators are applied in the same way as in Section~\ref{subsec:fermionic-excitation} already discussed.

The following algorithm describes the C++ function in its fundamental operation:
\begin{minted}[frame=lines, framesep=2mm, obeytabs=true, tabsize=4]{text}
Input: state_phi, state_psi, operator_sum (list of fermionic terms)

Initialize total_expectation_value = 0

For each term in operator_sum:
    For each basis_state in state_psi:

        Try applying all annihilation operators:
            If an orbital is already empty, skip to next basis_state
            Otherwise, track fermionic phase and remove electron

        Try applying all creation operators:
            If an orbital is already full, skip to next basis_state
            Otherwise, track fermionic phase and add electron

        Generate resulting new_basis_state

        Search for new_basis_state in state_phi:
            If found:
                Calculate overlap contribution (using amplitudes, term weight, and phase)
                Add contribution to total_expectation_value

Return total_expectation_value
\end{minted}

Apart from the efficiency mechanisms already discussed in~\ref{subsec:fermionic-excitation} this function implements a constant-time hashmap lookup to find the new basis state.
The implemented C++ function can be found in the appendix under~\ref{subsec:apendix-expectation-value}.